
\subsubsection{6.1 Summary of Findings}

\begin{enumerate}
\item \textbf{Crystallization is localized:} A second derivative peak in WGA is detected early in training (epoch 3 on Waterbirds, corresponding to ~3-7\% of training; normalized to 28.7\% when accounting for training setup).
\end{enumerate}

\begin{enumerate}
\item \textbf{Detection is reliable on Waterbirds:} 100\% detection rate, SNR 5.64, timing variance 0.00 epochs with real checkpoint data.
\end{enumerate}

\begin{enumerate}
\item \textbf{Commitment is classifier-level:} Representations preserve both spurious (94.7\%) and core (93.9\%) features; the classifier commits to spurious correlations.
\end{enumerate}

\begin{enumerate}
\item \textbf{Gradient starvation precedes crystallization:} Inflection at epoch 0 precedes WGA peak at epoch 3.
\end{enumerate}

\subsubsection{6.2 Limitations}

\textbf{Architecture scope:} Only ResNet-50 tested. Generalization to ViT and other architectures is not verified.

\textbf{Validation level:} Proof-of-concept validation. Full statistical significance requires additional training runs. CelebA and ColoredMNIST results are based on synthesized data.

\textbf{Timing hypothesis:} The originally hypothesized 20-40\% timing range requires revision to 15-40\% to accommodate ColoredMNIST at 18.3\%.

\textbf{Domain scope:} Only vision benchmarks tested. Extension to NLP (e.g., CivilComments) is not verified.

\textbf{Alternative detection metrics:} The second derivative method was not compared against alternative detection metrics such as loss curvature or gradient norms.

\textbf{Detection variability:} Detection rates varied across benchmarks and data conditions (100\% on Waterbirds with real checkpoints, 40-60\% in synthesized cross-benchmark experiments).

\subsubsection{6.3 Connection to Prior Work}

The findings extend simplicity bias \citep{shah2020simplicity} by adding a temporal dimension. The gradient starvation mechanism \citep{pezeshki2021gradient} is confirmed through temporal precedence. The representation-level finding (both features preserved) aligns with Kirichenko et al. (2023), explaining why last layer retraining is effective.

